\section*{Приложение В. Заголовочный файл и реализация класса \texttt{Validator}, а также \texttt{main.cpp}}
\addcontentsline{toc}{section}{Приложение В. Заголовочный файл и реализация класса \texttt{Validator}, а также файл \texttt{main.cpp}}

\begin{lstlisting}[language=C++, numbers=left, basicstyle=\ttfamily\scriptsize]
//main.cpp
#include <QApplication>
#include "mainwindow.h"

int main(int argc, char *argv[]) {
    QApplication app(argc, argv);

    MainWindow window;
    window.show();

    return app.exec();
}


// validator.h
#pragma once

#include <QString>
#include <QRegularExpression>
#include <QDate>

class Validator {
public:
    struct Result {
        bool valid;
        QString error;
        QString cleaned;

        Result(bool v = false, const QString& e = "", const QString& c = "")
            : valid(v), error(e), cleaned(c) {}
    };

    static Result validateName(const QString& name);
    static Result validatePhone(const QString& phone);
    static Result validateEmail(const QString& email);
    static Result validateBirthDate(const QDate& date);
    static Result validateAddress(const QString& address);

private:
    static QString trimSpaces(const QString& str);
    static bool isLeapYear(int year);
    static int daysInMonth(int month, int year);
};

// validator.cpp
#include "validator.h"

QString Validator::trimSpaces(const QString& str) {
    return str.trimmed().simplified();
}

bool Validator::isLeapYear(int year) {
    return (year % 4 == 0 && year % 100 != 0) || (year % 400 == 0);
}

int Validator::daysInMonth(int month, int year) {
    switch (month) {
        case 1: case 3: case 5: case 7: case 8: case 10: case 12:
            return 31;
        case 4: case 6: case 9: case 11:
            return 30;
        case 2:
            return isLeapYear(year) ? 29 : 28;
        default:
            return 0;
    }
}

Validator::Result Validator::validateName(const QString& name) {
    QString cleaned = trimSpaces(name);

    if (cleaned.isEmpty()) {
        return Result(false, "Поле не может быть пустым", "");
    }

    // - Фамилия, Имя или Отчество --- должны содержать только буквы
    // и цифры различных алфавитов, а также дефис и пробел, но при этом
    // должны начинаться только на заглавные буквы, и не могли бы
    // оканчиваться или начинаться на дефис. Все незначимые пробелы
    // перед и после данных должны удаляться.

    QRegularExpression startsWithUpper("^\\p{Lu}");
    if (!startsWithUpper.match(cleaned).hasMatch()) {
        return Result(false, "Должно начинаться с заглавной буквы", cleaned);
    }

    if (cleaned.startsWith('-') || cleaned.endsWith('-')) {
        return Result(false, "Не может начинаться или заканчиваться на дефис", cleaned);
    }

    QRegularExpression validChars("^[\\p{L}\\p{N}\\- ]+$");
    if (!validChars.match(cleaned).hasMatch()) {
        return Result(false, "Допустимы только буквы, цифры, дефис и пробел", cleaned);
    }

    return Result(true, "", cleaned);
}

Validator::Result Validator::validatePhone(const QString& phone) {
    QString cleaned = trimSpaces(phone);

    if (cleaned.isEmpty()) {
        return Result(false, "Телефон не может быть пустым", "");
    }

    // - Телефон должен быть записан в международном формате,
    // а храниться внутри как последовательность цифр.
    // Один из вариантов записи телефона:
    // префикс_выхода_на международную_линию код_страны код_города телефон.
    // Пример, +7 812 1234567, 8(800)123-1212, + 38 (032) 123 11 11

    QRegularExpression validFormat("^[\\+]?[\\d\\s\\(\\)\\-]+$");
    if (!validFormat.match(cleaned).hasMatch()) {
        return Result(false, "Недопустимые символы в номере телефона", cleaned);
    }

    QString digitsOnly;
    for (const QChar& c : cleaned) {
        if (c.isDigit()) {
            digitsOnly += c;
        }
    }

    if (digitsOnly.length() < 7) {
        return Result(false, "Телефон должен содержать минимум 7 цифр", digitsOnly);
    }

    if (digitsOnly.length() > 15) {
        return Result(false, "Телефон слишком длинный", digitsOnly);
    }

    return Result(true, "", digitsOnly);
}

Validator::Result Validator::validateEmail(const QString& email) {
    QString cleaned = email;
    cleaned.remove(' ');
    cleaned = cleaned.trimmed();

    if (cleaned.isEmpty()) {
        return Result(false, "Email не может быть пустым", "");
    }

    // - E-mail должен содержать в себе имя пользователя состоящее
    //  из латинских букв и цифр, символ разделения пользователя
    //  и имени домена(@), а также сам домен состоящий из латинских букв и цифр.
    //  Все незначащие пробелы (включая пробелы перед и после символа @)
    //  должны быть удалены.

    QRegularExpression emailRegex(
        "^[a-zA-Z0-9._-]+@[a-zA-Z0-9.-]+\\.[a-zA-Z]{2,}$"
    );

    if (!emailRegex.match(cleaned).hasMatch()) {
        return Result(false, "Неверный формат email", cleaned);
    }

    return Result(true, "", cleaned);
}

Validator::Result Validator::validateBirthDate(const QDate& date) {
    if (!date.isValid()) {
        return Result(false, "Некорректная дата", "");
    }

    // - Дата рождения должна быть меньше текущей даты, число месяцев в
    // дате должно быть от 1 до 12, число дней от 1 до 31, причем должно
    // учитываться различное число дней в месяце и високосные года.
    // Для обработки даты как вариант можно использовать классы
    // QDateEdit или QCalendarWidget.

    QDate today = QDate::currentDate();

    if (date >= today) {
        return Result(false, "Дата рождения должна быть меньше текущей даты", "");
    }

    int year = date.year();
    int month = date.month();
    int day = date.day();

    if (month < 1 || month > 12) {
        return Result(false, "Месяц должен быть от 1 до 12", "");
    }

    int maxDays = daysInMonth(month, year);
    if (day < 1 || day > maxDays) {
        return Result(false,
            QString("В этом месяце всего %1 дней").arg(maxDays), "");
    }

    return Result(true, "", date.toString("yyyy-MM-dd"));
}

Validator::Result Validator::validateAddress(const QString& address) {
    QString cleaned = trimSpaces(address);

    if (cleaned.isEmpty()) {
        return Result(false, "Адрес не может быть пустым", "");
    }

    return Result(true, "", cleaned);
}

\end{lstlisting}